% ============================================================================
%  CAPÍTULO II (continuación): ALCANCE Y LÍMITES
% ============================================================================

\section{Alcance}

\subsection{Tecnologías y entorno}

El proyecto opera sobre el siguiente stack tecnológico:

\begin{itemize}
  \item \textbf{Motor de base de datos:} PostgreSQL 15 hospedado en Supabase (región \textit{us-west-2}, pool de conexiones PgBouncer en modo \textit{session}).
  \item \textbf{ORM:} Entity Framework Core 10 con proveedor Npgsql 9 para .NET 10.
  \item \textbf{Aplicación web:} ASP.NET Core MVC 10 con ASP.NET Identity para autenticación y autorización basada en roles.
  \item \textbf{Lenguaje del servidor:} C\# 13.
  \item \textbf{Control de versiones:} Git / GitHub (\texttt{proyecto-final-bd-2}).
\end{itemize}

\subsection{Cobertura del rediseño}

El alcance del rediseño de base de datos comprende:

\begin{enumerate}
  \item Separación del esquema monolítico \texttt{public} en cinco esquemas especializados: \texttt{catalogo}, \texttt{ventas}, \texttt{operaciones}, \texttt{seguridad} y \texttt{auditoria}.
  \item Normalización de las tablas de negocio existentes hasta la 3FN.
  \item Creación de tablas nuevas: \texttt{Categorias}, \texttt{Marcas}, \texttt{Regiones}, \texttt{AuditoriaPedidos}, \texttt{AuditoriaComponentes}.
  \item Adición de columnas: \texttt{ItemsJson} migrado de \texttt{text} a \texttt{jsonb} en \texttt{CarritosAbandonados}.
  \item Definición de restricciones \texttt{CHECK} y \texttt{UNIQUE} en todas las tablas de negocio.
  \item Script SQL completo (\texttt{anexos/script\_bd2\_completo.sql}) con DDL, DML y objetos avanzados.
  \item Integración de la nueva estructura en EF Core mediante Fluent API (documentada en este informe, sin modificar el código fuente durante la elaboración del mismo).
\end{enumerate}

\subsection{Límites}

Los siguientes aspectos quedan fuera del alcance del presente trabajo:

\begin{itemize}
  \item Migración efectiva de la base de datos de Supabase en producción (el script es ejecutable pero no se aplicó al entorno productivo durante la elaboración del informe).
  \item Rediseño de la interfaz de usuario o de la lógica de negocio en C\#.
  \item Configuración de alta disponibilidad, replicación o particionamiento de tablas.
  \item Implementación de Row Level Security (RLS) de Supabase (se menciona como recomendación futura).
\end{itemize}
